% Standalone draft of the novelty section. Not yet \input into paper.tex.
% Uses only labels/cite keys already defined in paper.tex. No results here:
% every "how well" claim is deferred to Section~\ref{sec:evaluation}.
%
% To include: % Standalone draft of the novelty section. Not yet \input into paper.tex.
% Uses only labels/cite keys already defined in paper.tex. No results here:
% every "how well" claim is deferred to Section~\ref{sec:evaluation}.
%
% To include: % Standalone draft of the novelty section. Not yet \input into paper.tex.
% Uses only labels/cite keys already defined in paper.tex. No results here:
% every "how well" claim is deferred to Section~\ref{sec:evaluation}.
%
% To include: % Standalone draft of the novelty section. Not yet \input into paper.tex.
% Uses only labels/cite keys already defined in paper.tex. No results here:
% every "how well" claim is deferred to Section~\ref{sec:evaluation}.
%
% To include: \input{novelty} after the Related Work section.

\section{Novelty and Positioning}
\label{sec:novelty}

Each building block of our pipeline (sentence embeddings, an NLI
cross-encoder, Allen's algebra, finite-state machines) is individually
established. The contribution is \emph{what is decided with them, and in
what order}: a news claim is not judged by one score, but passes through a
sequence of gates in which every stage can only \emph{demote} the previous
stage's verdict for a stated, inspectable reason. This section states what is
new, how it works, and how it differs from the prior work of
Section~\ref{sec:related}.

\subsection{The Gap}

A system that reports ``is this claim about China--Romania relations
credible?'' over a live news stream faces three problems that existing
resources address only separately:

\begin{enumerate}
  \item \textbf{Same event, different words.} Independent outlets describe
        one event in different language, so exact-signature grouping
        under-counts corroboration.
  \item \textbf{Disagreement is not contradiction.} Two reports can
        disagree because the story \emph{moved on}
        (\emph{proposed} $\to$ \emph{signed}), not because one is wrong. A
        stateless contradiction classifier cannot tell the two apart.
  \item \textbf{Many domains are not many sources.} State media republish
        one wire story across language editions; counting domains rewards
        mirroring.
\end{enumerate}

\subsection{Differences from Prior Work}

\begin{table*}[t]
\caption{Positioning against the prior work discussed in
Section~\ref{sec:related}. Entries describe what each line of work does as
characterized there; ``---'' means the cited work does not address it.}
\label{tab:positioning}
\centering
\small
\begin{tabular}{p{2.6cm}p{3.2cm}p{3.4cm}p{3.4cm}p{3.4cm}}
\toprule
 & \textbf{GDELT / CAMEO}~\cite{leetaru2013gdelt,gerner2002cameo}
 & \textbf{General TKGs}~\cite{gottschalk2019eventkg,jain2020timeplex}
 & \textbf{Implicit temporal fact-checking}~\cite{allein2023implicit}
 & \textbf{This work} \\
\midrule
Domain
 & Global, all dyads
 & General-purpose
 & Claim vs.\ evidence pairs
 & One dyad, multilingual, full article text \\
Interval relation
 & ---
 & Validity interval per fact
 & Learned bucketed distance in days
 & Exact 13-relation Allen algebra between events \\
Corroboration
 & Mention counts
 & ---
 & ---
 & Semantic merge; publisher-network-aware diversity \\
Contradiction
 & ---
 & ---
 & Neural verdict, time as a feature
 & NLI score \emph{gated} by Allen relation and lifecycle \\
Process knowledge
 & ---
 & ---
 & ---
 & Per-event-type lifecycle state machines \\
Inference-time dependencies
 & ---
 & ---
 & Trained verification model
 & Pretrained models only; no training, no external fact-check API \\
\bottomrule
\end{tabular}
\end{table*}

The sharpest contrast is with Allein \textit{et al.}~\cite{allein2023implicit}.
They give a neural verifier a \emph{learned embedding of a time bucket} so it
can use temporal distance implicitly. We instead compute the \emph{exact
qualitative relation} between two intervals and use it as a hard, readable
rule. The trade-off is deliberate: a learned feature can capture patterns we
did not anticipate, but it cannot say \emph{why} a claim was or was not
treated as contradictory, and it needs labeled training data we do not have
for this dyad. Section~\ref{sec:evaluation} tests the two mechanisms against
each other directly.

\subsection{How the Gated Verdict Works}

Let a corroboration group $G$ contain claims $c_1,\dots,c_n$, each with a
sentence $s_i$, an event type $t_i$, and a date interval $I_i$ (possibly
unresolved). Verification proceeds in four stages; each stage may only
\emph{refine} the previous verdict.

\paragraph{Stage 1: semantic merge.}
Two exact-signature groups of the same event type are merged when
$\cos(e(s_a), e(s_b)) \ge \tau_m$, where $e(\cdot)$ is the mean-pooled
MiniLM embedding~\cite{wang2020minilm} and $\tau_m{=}0.75$. Merging is
transitive (union-find), and is restricted to a single event type so that two
different kinds of event are never fused merely for sounding alike.

\paragraph{Stage 2: NLI disagreement.}
For distinct sentences $s_a, s_b \in G$ with
$\cos(e(s_a),e(s_b)) \ge \tau_r$ ($\tau_r{=}0.15$), the contradiction score is
\begin{equation}
  d(a,b) = \max\bigl(P_{\text{NLI}}(\text{contra}\mid s_a, s_b),\;
                     P_{\text{NLI}}(\text{contra}\mid s_b, s_a)\bigr),
\end{equation}
and the pair is a \emph{candidate} disagreement if $d(a,b) \ge \tau_n$
($\tau_n{=}0.5$). NLI is directional, and which of two independently
published reports is the premise is arbitrary, so both orders are scored.
The relevance floor $\tau_r$ exists because NLI models trained on MNLI-style
data label many off-topic pairs \emph{contradiction} where a person would say
\emph{neutral}.

\paragraph{Stage 3: Allen gate.}
Each candidate is classified by the Allen relation
$R(I_a, I_b) \in \mathcal{A}$, $|\mathcal{A}|{=}13$:
\begin{equation}
  \text{verdict}(a,b) =
  \begin{cases}
    \textsc{evolution}      & R \in \{\,<,\ >,\ m,\ mi\,\},\\
    \textsc{contradiction}  & R \in \mathcal{A}\setminus\{<,>,m,mi\},\\
    \textsc{contradiction}  & I_a \text{ or } I_b \text{ unresolved.}
  \end{cases}
\end{equation}
The intuition is that two disagreeing claims about \emph{non-coincident}
moments describe a changing world, whereas disagreement about an
\emph{overlapping} moment describes a conflict. The last case is conservative
by design: absent temporal evidence we never reclassify.

\paragraph{Stage 4: lifecycle verification.}
The pairwise gate accepts \emph{any} sequential disagreement as evolution, so
it cannot see that ``ratified'' reported \emph{before} ``signed'' is not a
plausible story. For each event type $t$ we define a directed graph
$L_t=(S_t, E_t)$ of states and legal transitions
(Table~\ref{tab:lifecycles}). A sentence is assigned the state
\begin{equation}
  \hat{q}(s) = \arg\max_{q\in S_t}\; P_{\text{NLI}}(\text{entail}\mid s,\ h_q)
  \quad\text{if}\ \max \ge 0.5,
\end{equation}
where $h_q$ is a fixed natural-language template for state $q$ (zero-shot; no
training). Classified claims in a group flagged \textsc{evolution} are sorted
by date, and the sequence is checked against $E_t$. A consecutive pair
$(q_i,q_{i+1})\notin E_t$ with $q_i\ne q_{i+1}$ overrides the verdict to
\textsc{lifecycle\_violation}; repeats ($q_i=q_{i+1}$, e.g.\ two sources both
reporting \emph{signed}) are never violations.

\begin{algorithm}[t]
\caption{Gated claim verdict for one group (sketch)}
\abel{fig:gate-alg}
\begin{algorithmic}[1]
\STATE $G \leftarrow$ \textsc{SemanticMerge}$(G)$
\FORALL{pairs $(a,b)$ in $G$ with $\cos \ge \tau_r$}
  \IF{$d(a,b) \ge \tau_n$}
    \STATE $v \leftarrow$ \textsc{AllenGate}$(I_a, I_b)$
    \STATE record $(a,b,v)$
  \ENDIF
\ENDFOR
\IF{any recorded $v=\textsc{evolution}$}
  \STATE $\sigma \leftarrow$ date-sorted states $\hat q(s_i)$, dropping unclassified
  \IF{$\exists i:\ (\sigma_i,\sigma_{i+1})\notin E_t \wedge \sigma_i\ne\sigma_{i+1}$}
    \STATE override those $v$ to \textsc{lifecycle\_violation}
  \ENDIF
\ENDIF
\RETURN verdicts, with \texttt{source\_diversity} counted over publisher networks
\end{algorithmic}
\end{figure}

\subsection{Publisher-Network-Aware Diversity}

Independence is measured over \emph{publisher networks}, not domains:
$\text{div}(G)=\bigl|\{\nu(\text{dom}(c)) : c\in G\}\bigr|$, where $\nu$ maps
a domain to a configured network identifier and is the identity on unmapped
domains. Because $\nu$ is a function, $\text{div}$ can only \emph{decrease}
relative to raw domain counting, and only for known multi-domain operations;
it cannot manufacture diversity.

\subsection{Why This Is Different, Not Merely Larger}

\begin{itemize}
  \item \textbf{Auditable by construction.} Every demotion has a named cause
        (below relevance floor, non-coincident interval, illegal
        transition), so a flagged claim carries a reason, not only a score.
  \item \textbf{No training at inference.} All models are pretrained and used
        zero-shot; the system needs no labeled data for the dyad to run.
  \item \textbf{Failure modes are measured, not assumed.} We evaluate each
        stage in isolation on real mined pairs and report where a stage
        \emph{fails} (Section~\ref{sec:evaluation}), including a stage-level
        precision collapse of raw NLI on formulaic press-release text that a
        hand-authored test set did not reveal.
  \item \textbf{Fidelity to the domain.} Full article text (not headlines) is
        ingested; machine-translation failures are detected rather than
        silently propagated; and out-of-scope matches such as sports fixtures
        are excluded before extraction.
\end{itemize}

\subsection{Scope of the Claim}

We do \emph{not} claim a new NLI model, a new embedding model, or a new
temporal algebra. We claim (i) the composition above and its ordering,
(ii) a measured comparison of exact interval reasoning against a bucketed
heuristic in the same pipeline, and (iii) a documented set of failure modes of
off-the-shelf NLI on real bilateral news. The lifecycle models are
hand-authored and not yet validated against the diplomatic or economic
literature; we treat them as a first-pass domain ontology and say so
(Section~\ref{sec:discussion}).
 after the Related Work section.

\section{Novelty and Positioning}
\label{sec:novelty}

Each building block of our pipeline (sentence embeddings, an NLI
cross-encoder, Allen's algebra, finite-state machines) is individually
established. The contribution is \emph{what is decided with them, and in
what order}: a news claim is not judged by one score, but passes through a
sequence of gates in which every stage can only \emph{demote} the previous
stage's verdict for a stated, inspectable reason. This section states what is
new, how it works, and how it differs from the prior work of
Section~\ref{sec:related}.

\subsection{The Gap}

A system that reports ``is this claim about China--Romania relations
credible?'' over a live news stream faces three problems that existing
resources address only separately:

\begin{enumerate}
  \item \textbf{Same event, different words.} Independent outlets describe
        one event in different language, so exact-signature grouping
        under-counts corroboration.
  \item \textbf{Disagreement is not contradiction.} Two reports can
        disagree because the story \emph{moved on}
        (\emph{proposed} $\to$ \emph{signed}), not because one is wrong. A
        stateless contradiction classifier cannot tell the two apart.
  \item \textbf{Many domains are not many sources.} State media republish
        one wire story across language editions; counting domains rewards
        mirroring.
\end{enumerate}

\subsection{Differences from Prior Work}

\begin{table*}[t]
\caption{Positioning against the prior work discussed in
Section~\ref{sec:related}. Entries describe what each line of work does as
characterized there; ``---'' means the cited work does not address it.}
\label{tab:positioning}
\centering
\small
\begin{tabular}{p{2.6cm}p{3.2cm}p{3.4cm}p{3.4cm}p{3.4cm}}
\toprule
 & \textbf{GDELT / CAMEO}~\cite{leetaru2013gdelt,gerner2002cameo}
 & \textbf{General TKGs}~\cite{gottschalk2019eventkg,jain2020timeplex}
 & \textbf{Implicit temporal fact-checking}~\cite{allein2023implicit}
 & \textbf{This work} \\
\midrule
Domain
 & Global, all dyads
 & General-purpose
 & Claim vs.\ evidence pairs
 & One dyad, multilingual, full article text \\
Interval relation
 & ---
 & Validity interval per fact
 & Learned bucketed distance in days
 & Exact 13-relation Allen algebra between events \\
Corroboration
 & Mention counts
 & ---
 & ---
 & Semantic merge; publisher-network-aware diversity \\
Contradiction
 & ---
 & ---
 & Neural verdict, time as a feature
 & NLI score \emph{gated} by Allen relation and lifecycle \\
Process knowledge
 & ---
 & ---
 & ---
 & Per-event-type lifecycle state machines \\
Inference-time dependencies
 & ---
 & ---
 & Trained verification model
 & Pretrained models only; no training, no external fact-check API \\
\bottomrule
\end{tabular}
\end{table*}

The sharpest contrast is with Allein \textit{et al.}~\cite{allein2023implicit}.
They give a neural verifier a \emph{learned embedding of a time bucket} so it
can use temporal distance implicitly. We instead compute the \emph{exact
qualitative relation} between two intervals and use it as a hard, readable
rule. The trade-off is deliberate: a learned feature can capture patterns we
did not anticipate, but it cannot say \emph{why} a claim was or was not
treated as contradictory, and it needs labeled training data we do not have
for this dyad. Section~\ref{sec:evaluation} tests the two mechanisms against
each other directly.

\subsection{How the Gated Verdict Works}

Let a corroboration group $G$ contain claims $c_1,\dots,c_n$, each with a
sentence $s_i$, an event type $t_i$, and a date interval $I_i$ (possibly
unresolved). Verification proceeds in four stages; each stage may only
\emph{refine} the previous verdict.

\paragraph{Stage 1: semantic merge.}
Two exact-signature groups of the same event type are merged when
$\cos(e(s_a), e(s_b)) \ge \tau_m$, where $e(\cdot)$ is the mean-pooled
MiniLM embedding~\cite{wang2020minilm} and $\tau_m{=}0.75$. Merging is
transitive (union-find), and is restricted to a single event type so that two
different kinds of event are never fused merely for sounding alike.

\paragraph{Stage 2: NLI disagreement.}
For distinct sentences $s_a, s_b \in G$ with
$\cos(e(s_a),e(s_b)) \ge \tau_r$ ($\tau_r{=}0.15$), the contradiction score is
\begin{equation}
  d(a,b) = \max\bigl(P_{\text{NLI}}(\text{contra}\mid s_a, s_b),\;
                     P_{\text{NLI}}(\text{contra}\mid s_b, s_a)\bigr),
\end{equation}
and the pair is a \emph{candidate} disagreement if $d(a,b) \ge \tau_n$
($\tau_n{=}0.5$). NLI is directional, and which of two independently
published reports is the premise is arbitrary, so both orders are scored.
The relevance floor $\tau_r$ exists because NLI models trained on MNLI-style
data label many off-topic pairs \emph{contradiction} where a person would say
\emph{neutral}.

\paragraph{Stage 3: Allen gate.}
Each candidate is classified by the Allen relation
$R(I_a, I_b) \in \mathcal{A}$, $|\mathcal{A}|{=}13$:
\begin{equation}
  \text{verdict}(a,b) =
  \begin{cases}
    \textsc{evolution}      & R \in \{\,<,\ >,\ m,\ mi\,\},\\
    \textsc{contradiction}  & R \in \mathcal{A}\setminus\{<,>,m,mi\},\\
    \textsc{contradiction}  & I_a \text{ or } I_b \text{ unresolved.}
  \end{cases}
\end{equation}
The intuition is that two disagreeing claims about \emph{non-coincident}
moments describe a changing world, whereas disagreement about an
\emph{overlapping} moment describes a conflict. The last case is conservative
by design: absent temporal evidence we never reclassify.

\paragraph{Stage 4: lifecycle verification.}
The pairwise gate accepts \emph{any} sequential disagreement as evolution, so
it cannot see that ``ratified'' reported \emph{before} ``signed'' is not a
plausible story. For each event type $t$ we define a directed graph
$L_t=(S_t, E_t)$ of states and legal transitions
(Table~\ref{tab:lifecycles}). A sentence is assigned the state
\begin{equation}
  \hat{q}(s) = \arg\max_{q\in S_t}\; P_{\text{NLI}}(\text{entail}\mid s,\ h_q)
  \quad\text{if}\ \max \ge 0.5,
\end{equation}
where $h_q$ is a fixed natural-language template for state $q$ (zero-shot; no
training). Classified claims in a group flagged \textsc{evolution} are sorted
by date, and the sequence is checked against $E_t$. A consecutive pair
$(q_i,q_{i+1})\notin E_t$ with $q_i\ne q_{i+1}$ overrides the verdict to
\textsc{lifecycle\_violation}; repeats ($q_i=q_{i+1}$, e.g.\ two sources both
reporting \emph{signed}) are never violations.

\begin{algorithm}[t]
\caption{Gated claim verdict for one group (sketch)}
\abel{fig:gate-alg}
\begin{algorithmic}[1]
\STATE $G \leftarrow$ \textsc{SemanticMerge}$(G)$
\FORALL{pairs $(a,b)$ in $G$ with $\cos \ge \tau_r$}
  \IF{$d(a,b) \ge \tau_n$}
    \STATE $v \leftarrow$ \textsc{AllenGate}$(I_a, I_b)$
    \STATE record $(a,b,v)$
  \ENDIF
\ENDFOR
\IF{any recorded $v=\textsc{evolution}$}
  \STATE $\sigma \leftarrow$ date-sorted states $\hat q(s_i)$, dropping unclassified
  \IF{$\exists i:\ (\sigma_i,\sigma_{i+1})\notin E_t \wedge \sigma_i\ne\sigma_{i+1}$}
    \STATE override those $v$ to \textsc{lifecycle\_violation}
  \ENDIF
\ENDIF
\RETURN verdicts, with \texttt{source\_diversity} counted over publisher networks
\end{algorithmic}
\end{figure}

\subsection{Publisher-Network-Aware Diversity}

Independence is measured over \emph{publisher networks}, not domains:
$\text{div}(G)=\bigl|\{\nu(\text{dom}(c)) : c\in G\}\bigr|$, where $\nu$ maps
a domain to a configured network identifier and is the identity on unmapped
domains. Because $\nu$ is a function, $\text{div}$ can only \emph{decrease}
relative to raw domain counting, and only for known multi-domain operations;
it cannot manufacture diversity.

\subsection{Why This Is Different, Not Merely Larger}

\begin{itemize}
  \item \textbf{Auditable by construction.} Every demotion has a named cause
        (below relevance floor, non-coincident interval, illegal
        transition), so a flagged claim carries a reason, not only a score.
  \item \textbf{No training at inference.} All models are pretrained and used
        zero-shot; the system needs no labeled data for the dyad to run.
  \item \textbf{Failure modes are measured, not assumed.} We evaluate each
        stage in isolation on real mined pairs and report where a stage
        \emph{fails} (Section~\ref{sec:evaluation}), including a stage-level
        precision collapse of raw NLI on formulaic press-release text that a
        hand-authored test set did not reveal.
  \item \textbf{Fidelity to the domain.} Full article text (not headlines) is
        ingested; machine-translation failures are detected rather than
        silently propagated; and out-of-scope matches such as sports fixtures
        are excluded before extraction.
\end{itemize}

\subsection{Scope of the Claim}

We do \emph{not} claim a new NLI model, a new embedding model, or a new
temporal algebra. We claim (i) the composition above and its ordering,
(ii) a measured comparison of exact interval reasoning against a bucketed
heuristic in the same pipeline, and (iii) a documented set of failure modes of
off-the-shelf NLI on real bilateral news. The lifecycle models are
hand-authored and not yet validated against the diplomatic or economic
literature; we treat them as a first-pass domain ontology and say so
(Section~\ref{sec:discussion}).
 after the Related Work section.

\section{Novelty and Positioning}
\label{sec:novelty}

Each building block of our pipeline (sentence embeddings, an NLI
cross-encoder, Allen's algebra, finite-state machines) is individually
established. The contribution is \emph{what is decided with them, and in
what order}: a news claim is not judged by one score, but passes through a
sequence of gates in which every stage can only \emph{demote} the previous
stage's verdict for a stated, inspectable reason. This section states what is
new, how it works, and how it differs from the prior work of
Section~\ref{sec:related}.

\subsection{The Gap}

A system that reports ``is this claim about China--Romania relations
credible?'' over a live news stream faces three problems that existing
resources address only separately:

\begin{enumerate}
  \item \textbf{Same event, different words.} Independent outlets describe
        one event in different language, so exact-signature grouping
        under-counts corroboration.
  \item \textbf{Disagreement is not contradiction.} Two reports can
        disagree because the story \emph{moved on}
        (\emph{proposed} $\to$ \emph{signed}), not because one is wrong. A
        stateless contradiction classifier cannot tell the two apart.
  \item \textbf{Many domains are not many sources.} State media republish
        one wire story across language editions; counting domains rewards
        mirroring.
\end{enumerate}

\subsection{Differences from Prior Work}

\begin{table*}[t]
\caption{Positioning against the prior work discussed in
Section~\ref{sec:related}. Entries describe what each line of work does as
characterized there; ``---'' means the cited work does not address it.}
\label{tab:positioning}
\centering
\small
\begin{tabular}{p{2.6cm}p{3.2cm}p{3.4cm}p{3.4cm}p{3.4cm}}
\toprule
 & \textbf{GDELT / CAMEO}~\cite{leetaru2013gdelt,gerner2002cameo}
 & \textbf{General TKGs}~\cite{gottschalk2019eventkg,jain2020timeplex}
 & \textbf{Implicit temporal fact-checking}~\cite{allein2023implicit}
 & \textbf{This work} \\
\midrule
Domain
 & Global, all dyads
 & General-purpose
 & Claim vs.\ evidence pairs
 & One dyad, multilingual, full article text \\
Interval relation
 & ---
 & Validity interval per fact
 & Learned bucketed distance in days
 & Exact 13-relation Allen algebra between events \\
Corroboration
 & Mention counts
 & ---
 & ---
 & Semantic merge; publisher-network-aware diversity \\
Contradiction
 & ---
 & ---
 & Neural verdict, time as a feature
 & NLI score \emph{gated} by Allen relation and lifecycle \\
Process knowledge
 & ---
 & ---
 & ---
 & Per-event-type lifecycle state machines \\
Inference-time dependencies
 & ---
 & ---
 & Trained verification model
 & Pretrained models only; no training, no external fact-check API \\
\bottomrule
\end{tabular}
\end{table*}

The sharpest contrast is with Allein \textit{et al.}~\cite{allein2023implicit}.
They give a neural verifier a \emph{learned embedding of a time bucket} so it
can use temporal distance implicitly. We instead compute the \emph{exact
qualitative relation} between two intervals and use it as a hard, readable
rule. The trade-off is deliberate: a learned feature can capture patterns we
did not anticipate, but it cannot say \emph{why} a claim was or was not
treated as contradictory, and it needs labeled training data we do not have
for this dyad. Section~\ref{sec:evaluation} tests the two mechanisms against
each other directly.

\subsection{How the Gated Verdict Works}

Let a corroboration group $G$ contain claims $c_1,\dots,c_n$, each with a
sentence $s_i$, an event type $t_i$, and a date interval $I_i$ (possibly
unresolved). Verification proceeds in four stages; each stage may only
\emph{refine} the previous verdict.

\paragraph{Stage 1: semantic merge.}
Two exact-signature groups of the same event type are merged when
$\cos(e(s_a), e(s_b)) \ge \tau_m$, where $e(\cdot)$ is the mean-pooled
MiniLM embedding~\cite{wang2020minilm} and $\tau_m{=}0.75$. Merging is
transitive (union-find), and is restricted to a single event type so that two
different kinds of event are never fused merely for sounding alike.

\paragraph{Stage 2: NLI disagreement.}
For distinct sentences $s_a, s_b \in G$ with
$\cos(e(s_a),e(s_b)) \ge \tau_r$ ($\tau_r{=}0.15$), the contradiction score is
\begin{equation}
  d(a,b) = \max\bigl(P_{\text{NLI}}(\text{contra}\mid s_a, s_b),\;
                     P_{\text{NLI}}(\text{contra}\mid s_b, s_a)\bigr),
\end{equation}
and the pair is a \emph{candidate} disagreement if $d(a,b) \ge \tau_n$
($\tau_n{=}0.5$). NLI is directional, and which of two independently
published reports is the premise is arbitrary, so both orders are scored.
The relevance floor $\tau_r$ exists because NLI models trained on MNLI-style
data label many off-topic pairs \emph{contradiction} where a person would say
\emph{neutral}.

\paragraph{Stage 3: Allen gate.}
Each candidate is classified by the Allen relation
$R(I_a, I_b) \in \mathcal{A}$, $|\mathcal{A}|{=}13$:
\begin{equation}
  \text{verdict}(a,b) =
  \begin{cases}
    \textsc{evolution}      & R \in \{\,<,\ >,\ m,\ mi\,\},\\
    \textsc{contradiction}  & R \in \mathcal{A}\setminus\{<,>,m,mi\},\\
    \textsc{contradiction}  & I_a \text{ or } I_b \text{ unresolved.}
  \end{cases}
\end{equation}
The intuition is that two disagreeing claims about \emph{non-coincident}
moments describe a changing world, whereas disagreement about an
\emph{overlapping} moment describes a conflict. The last case is conservative
by design: absent temporal evidence we never reclassify.

\paragraph{Stage 4: lifecycle verification.}
The pairwise gate accepts \emph{any} sequential disagreement as evolution, so
it cannot see that ``ratified'' reported \emph{before} ``signed'' is not a
plausible story. For each event type $t$ we define a directed graph
$L_t=(S_t, E_t)$ of states and legal transitions
(Table~\ref{tab:lifecycles}). A sentence is assigned the state
\begin{equation}
  \hat{q}(s) = \arg\max_{q\in S_t}\; P_{\text{NLI}}(\text{entail}\mid s,\ h_q)
  \quad\text{if}\ \max \ge 0.5,
\end{equation}
where $h_q$ is a fixed natural-language template for state $q$ (zero-shot; no
training). Classified claims in a group flagged \textsc{evolution} are sorted
by date, and the sequence is checked against $E_t$. A consecutive pair
$(q_i,q_{i+1})\notin E_t$ with $q_i\ne q_{i+1}$ overrides the verdict to
\textsc{lifecycle\_violation}; repeats ($q_i=q_{i+1}$, e.g.\ two sources both
reporting \emph{signed}) are never violations.

\begin{algorithm}[t]
\caption{Gated claim verdict for one group (sketch)}
\abel{fig:gate-alg}
\begin{algorithmic}[1]
\STATE $G \leftarrow$ \textsc{SemanticMerge}$(G)$
\FORALL{pairs $(a,b)$ in $G$ with $\cos \ge \tau_r$}
  \IF{$d(a,b) \ge \tau_n$}
    \STATE $v \leftarrow$ \textsc{AllenGate}$(I_a, I_b)$
    \STATE record $(a,b,v)$
  \ENDIF
\ENDFOR
\IF{any recorded $v=\textsc{evolution}$}
  \STATE $\sigma \leftarrow$ date-sorted states $\hat q(s_i)$, dropping unclassified
  \IF{$\exists i:\ (\sigma_i,\sigma_{i+1})\notin E_t \wedge \sigma_i\ne\sigma_{i+1}$}
    \STATE override those $v$ to \textsc{lifecycle\_violation}
  \ENDIF
\ENDIF
\RETURN verdicts, with \texttt{source\_diversity} counted over publisher networks
\end{algorithmic}
\end{figure}

\subsection{Publisher-Network-Aware Diversity}

Independence is measured over \emph{publisher networks}, not domains:
$\text{div}(G)=\bigl|\{\nu(\text{dom}(c)) : c\in G\}\bigr|$, where $\nu$ maps
a domain to a configured network identifier and is the identity on unmapped
domains. Because $\nu$ is a function, $\text{div}$ can only \emph{decrease}
relative to raw domain counting, and only for known multi-domain operations;
it cannot manufacture diversity.

\subsection{Why This Is Different, Not Merely Larger}

\begin{itemize}
  \item \textbf{Auditable by construction.} Every demotion has a named cause
        (below relevance floor, non-coincident interval, illegal
        transition), so a flagged claim carries a reason, not only a score.
  \item \textbf{No training at inference.} All models are pretrained and used
        zero-shot; the system needs no labeled data for the dyad to run.
  \item \textbf{Failure modes are measured, not assumed.} We evaluate each
        stage in isolation on real mined pairs and report where a stage
        \emph{fails} (Section~\ref{sec:evaluation}), including a stage-level
        precision collapse of raw NLI on formulaic press-release text that a
        hand-authored test set did not reveal.
  \item \textbf{Fidelity to the domain.} Full article text (not headlines) is
        ingested; machine-translation failures are detected rather than
        silently propagated; and out-of-scope matches such as sports fixtures
        are excluded before extraction.
\end{itemize}

\subsection{Scope of the Claim}

We do \emph{not} claim a new NLI model, a new embedding model, or a new
temporal algebra. We claim (i) the composition above and its ordering,
(ii) a measured comparison of exact interval reasoning against a bucketed
heuristic in the same pipeline, and (iii) a documented set of failure modes of
off-the-shelf NLI on real bilateral news. The lifecycle models are
hand-authored and not yet validated against the diplomatic or economic
literature; we treat them as a first-pass domain ontology and say so
(Section~\ref{sec:discussion}).
 after the Related Work section.

\section{Novelty and Positioning}
\label{sec:novelty}

Each building block of our pipeline (sentence embeddings, an NLI
cross-encoder, Allen's algebra, finite-state machines) is individually
established. The contribution is \emph{what is decided with them, and in
what order}: a news claim is not judged by one score, but passes through a
sequence of gates in which every stage can only \emph{demote} the previous
stage's verdict for a stated, inspectable reason. This section states what is
new, how it works, and how it differs from the prior work of
Section~\ref{sec:related}.

\subsection{The Gap}

A system that reports ``is this claim about China--Romania relations
credible?'' over a live news stream faces three problems that existing
resources address only separately:

\begin{enumerate}
  \item \textbf{Same event, different words.} Independent outlets describe
        one event in different language, so exact-signature grouping
        under-counts corroboration.
  \item \textbf{Disagreement is not contradiction.} Two reports can
        disagree because the story \emph{moved on}
        (\emph{proposed} $\to$ \emph{signed}), not because one is wrong. A
        stateless contradiction classifier cannot tell the two apart.
  \item \textbf{Many domains are not many sources.} State media republish
        one wire story across language editions; counting domains rewards
        mirroring.
\end{enumerate}

\subsection{Differences from Prior Work}

\begin{table*}[t]
\caption{Positioning against the prior work discussed in
Section~\ref{sec:related}. Entries describe what each line of work does as
characterized there; ``---'' means the cited work does not address it.}
\label{tab:positioning}
\centering
\small
\begin{tabular}{p{2.6cm}p{3.2cm}p{3.4cm}p{3.4cm}p{3.4cm}}
\toprule
 & \textbf{GDELT / CAMEO}~\cite{leetaru2013gdelt,gerner2002cameo}
 & \textbf{General TKGs}~\cite{gottschalk2019eventkg,jain2020timeplex}
 & \textbf{Implicit temporal fact-checking}~\cite{allein2023implicit}
 & \textbf{This work} \\
\midrule
Domain
 & Global, all dyads
 & General-purpose
 & Claim vs.\ evidence pairs
 & One dyad, multilingual, full article text \\
Interval relation
 & ---
 & Validity interval per fact
 & Learned bucketed distance in days
 & Exact 13-relation Allen algebra between events \\
Corroboration
 & Mention counts
 & ---
 & ---
 & Semantic merge; publisher-network-aware diversity \\
Contradiction
 & ---
 & ---
 & Neural verdict, time as a feature
 & NLI score \emph{gated} by Allen relation and lifecycle \\
Process knowledge
 & ---
 & ---
 & ---
 & Per-event-type lifecycle state machines \\
Inference-time dependencies
 & ---
 & ---
 & Trained verification model
 & Pretrained models only; no training, no external fact-check API \\
\bottomrule
\end{tabular}
\end{table*}

The sharpest contrast is with Allein \textit{et al.}~\cite{allein2023implicit}.
They give a neural verifier a \emph{learned embedding of a time bucket} so it
can use temporal distance implicitly. We instead compute the \emph{exact
qualitative relation} between two intervals and use it as a hard, readable
rule. The trade-off is deliberate: a learned feature can capture patterns we
did not anticipate, but it cannot say \emph{why} a claim was or was not
treated as contradictory, and it needs labeled training data we do not have
for this dyad. Section~\ref{sec:evaluation} tests the two mechanisms against
each other directly.

\subsection{How the Gated Verdict Works}

Let a corroboration group $G$ contain claims $c_1,\dots,c_n$, each with a
sentence $s_i$, an event type $t_i$, and a date interval $I_i$ (possibly
unresolved). Verification proceeds in four stages; each stage may only
\emph{refine} the previous verdict.

\paragraph{Stage 1: semantic merge.}
Two exact-signature groups of the same event type are merged when
$\cos(e(s_a), e(s_b)) \ge \tau_m$, where $e(\cdot)$ is the mean-pooled
MiniLM embedding~\cite{wang2020minilm} and $\tau_m{=}0.75$. Merging is
transitive (union-find), and is restricted to a single event type so that two
different kinds of event are never fused merely for sounding alike.

\paragraph{Stage 2: NLI disagreement.}
For distinct sentences $s_a, s_b \in G$ with
$\cos(e(s_a),e(s_b)) \ge \tau_r$ ($\tau_r{=}0.15$), the contradiction score is
\begin{equation}
  d(a,b) = \max\bigl(P_{\text{NLI}}(\text{contra}\mid s_a, s_b),\;
                     P_{\text{NLI}}(\text{contra}\mid s_b, s_a)\bigr),
\end{equation}
and the pair is a \emph{candidate} disagreement if $d(a,b) \ge \tau_n$
($\tau_n{=}0.5$). NLI is directional, and which of two independently
published reports is the premise is arbitrary, so both orders are scored.
The relevance floor $\tau_r$ exists because NLI models trained on MNLI-style
data label many off-topic pairs \emph{contradiction} where a person would say
\emph{neutral}.

\paragraph{Stage 3: Allen gate.}
Each candidate is classified by the Allen relation
$R(I_a, I_b) \in \mathcal{A}$, $|\mathcal{A}|{=}13$:
\begin{equation}
  \text{verdict}(a,b) =
  \begin{cases}
    \textsc{evolution}      & R \in \{\,<,\ >,\ m,\ mi\,\},\\
    \textsc{contradiction}  & R \in \mathcal{A}\setminus\{<,>,m,mi\},\\
    \textsc{contradiction}  & I_a \text{ or } I_b \text{ unresolved.}
  \end{cases}
\end{equation}
The intuition is that two disagreeing claims about \emph{non-coincident}
moments describe a changing world, whereas disagreement about an
\emph{overlapping} moment describes a conflict. The last case is conservative
by design: absent temporal evidence we never reclassify.

\paragraph{Stage 4: lifecycle verification.}
The pairwise gate accepts \emph{any} sequential disagreement as evolution, so
it cannot see that ``ratified'' reported \emph{before} ``signed'' is not a
plausible story. For each event type $t$ we define a directed graph
$L_t=(S_t, E_t)$ of states and legal transitions
(Table~\ref{tab:lifecycles}). A sentence is assigned the state
\begin{equation}
  \hat{q}(s) = \arg\max_{q\in S_t}\; P_{\text{NLI}}(\text{entail}\mid s,\ h_q)
  \quad\text{if}\ \max \ge 0.5,
\end{equation}
where $h_q$ is a fixed natural-language template for state $q$ (zero-shot; no
training). Classified claims in a group flagged \textsc{evolution} are sorted
by date, and the sequence is checked against $E_t$. A consecutive pair
$(q_i,q_{i+1})\notin E_t$ with $q_i\ne q_{i+1}$ overrides the verdict to
\textsc{lifecycle\_violation}; repeats ($q_i=q_{i+1}$, e.g.\ two sources both
reporting \emph{signed}) are never violations.

\begin{algorithm}[t]
\caption{Gated claim verdict for one group (sketch)}
\abel{fig:gate-alg}
\begin{algorithmic}[1]
\STATE $G \leftarrow$ \textsc{SemanticMerge}$(G)$
\FORALL{pairs $(a,b)$ in $G$ with $\cos \ge \tau_r$}
  \IF{$d(a,b) \ge \tau_n$}
    \STATE $v \leftarrow$ \textsc{AllenGate}$(I_a, I_b)$
    \STATE record $(a,b,v)$
  \ENDIF
\ENDFOR
\IF{any recorded $v=\textsc{evolution}$}
  \STATE $\sigma \leftarrow$ date-sorted states $\hat q(s_i)$, dropping unclassified
  \IF{$\exists i:\ (\sigma_i,\sigma_{i+1})\notin E_t \wedge \sigma_i\ne\sigma_{i+1}$}
    \STATE override those $v$ to \textsc{lifecycle\_violation}
  \ENDIF
\ENDIF
\RETURN verdicts, with \texttt{source\_diversity} counted over publisher networks
\end{algorithmic}
\end{figure}

\subsection{Publisher-Network-Aware Diversity}

Independence is measured over \emph{publisher networks}, not domains:
$\text{div}(G)=\bigl|\{\nu(\text{dom}(c)) : c\in G\}\bigr|$, where $\nu$ maps
a domain to a configured network identifier and is the identity on unmapped
domains. Because $\nu$ is a function, $\text{div}$ can only \emph{decrease}
relative to raw domain counting, and only for known multi-domain operations;
it cannot manufacture diversity.

\subsection{Why This Is Different, Not Merely Larger}

\begin{itemize}
  \item \textbf{Auditable by construction.} Every demotion has a named cause
        (below relevance floor, non-coincident interval, illegal
        transition), so a flagged claim carries a reason, not only a score.
  \item \textbf{No training at inference.} All models are pretrained and used
        zero-shot; the system needs no labeled data for the dyad to run.
  \item \textbf{Failure modes are measured, not assumed.} We evaluate each
        stage in isolation on real mined pairs and report where a stage
        \emph{fails} (Section~\ref{sec:evaluation}), including a stage-level
        precision collapse of raw NLI on formulaic press-release text that a
        hand-authored test set did not reveal.
  \item \textbf{Fidelity to the domain.} Full article text (not headlines) is
        ingested; machine-translation failures are detected rather than
        silently propagated; and out-of-scope matches such as sports fixtures
        are excluded before extraction.
\end{itemize}

\subsection{Scope of the Claim}

We do \emph{not} claim a new NLI model, a new embedding model, or a new
temporal algebra. We claim (i) the composition above and its ordering,
(ii) a measured comparison of exact interval reasoning against a bucketed
heuristic in the same pipeline, and (iii) a documented set of failure modes of
off-the-shelf NLI on real bilateral news. The lifecycle models are
hand-authored and not yet validated against the diplomatic or economic
literature; we treat them as a first-pass domain ontology and say so
(Section~\ref{sec:discussion}).
